\subsection{Permute \& Reject}

  \begin{figure}[h]
    \begin{center}
      \includegraphics[width=0.3\textwidth]{img/permute16.eps}
    \end{center}
    \caption{Polygon mit 16 Punkten (Permute \& Reject)}
    \label{fig:permute16}
  \end{figure}

  \emph{Permute \& Reject} ist der einfachste in der Software enthaltenene
  Algorithmus zur Erzeugung von simplen Polygonen, implementiert nach einer
  Beschreibung von Martin Held in~\cite{held98polygons}.

  \subsubsection{Algorithmus}

    Der Algorithmus basiert auf der zufälligen Erzeugung von $n$ Kanten zwischen
    $n$ gegebenen Punkten, sodass von jedem Punkt zwei Kanten ausgehen.
    Anschließend wird überprüft, ob das entstandene Polygon \enquote{simpel},
    d.h. überschneidungsfrei ist. Falls nicht, wird das Polygon verworfen und
    von vorne begonnen.

    \begin{code}[caption={Permute \& Reject},label=permutelisting,mathescape=true]
      retry:
      $P$ $\leftarrow$ given set of points
      $p$ $\leftarrow$ remove random point from $P$
      while ($P$ not empty)
        $q$ $\leftarrow$ remove random point from $P$
        draw edge from $p$ to $q$
        $p$ $\leftarrow$ $q$
      if resulting polygon is NOT simple
        goto retry
    \end{code}

  \subsubsection{Eigenschaften}

    Ausgehend von einer beliebigen Punktmenge kann Permute \& Reject sämtliche
    möglichen Polygone erzeugen. Jedoch lässt die enorme Laufzeit des
    Algorithmus eine praktische Verwendung für Polygone mit vielen Punkten
    nicht zu. Eine -- wenn auch nicht scharfe -- obere Schranke für die
    Laufzeit lässt sich mit $\bigO(n!)$ beziffern, entsprechend dem
    Durchprobieren jeglicher möglichen Kombination von Punkten. In unseren
    Experimenten konnten wir kein Polygon mit mehr als 28 Punkten erzeugen.
  